\subsection{Überführung in die Produktentwicklung}
\label{subsec:ueberfuehrung-produktentwicklung}

Nach der Generierung entscheidet das eigenständige Produktprojekt über
Fachmodelle, Nutzerdaten, Persistenzprovider, Assets, Schema und Migrationen,
Backup und Recovery, Sicherheitsmodell und gegebenenfalls Synchronisation
\parencite{kleiveist2026persistence}. Geschäftsregeln, Oberfläche,
Integrationen, Betrieb, Skalierung, Monitoring, Datenschutz und Signing bleiben
ebenfalls Produktaufgaben. Die Baseline strukturiert diese Arbeit, löst sie
aber nicht.

Master und Produkt besitzen getrennte Lebenszyklen. Es gibt keinen dauerhaften
Update- oder Synchronisationsmechanismus. Master-Änderungen müssen gezielt in
Produkte übernommen werden; Produktänderungen gehören nur nach
Generalisierungs- und Verifikationsprüfung zurück in die Baseline.

Template- und Produktabnahme bleiben getrennt. Die commitgebundene Abnahme
belegt Generierung, Installation, Tests und ausgewählte Builds der Baseline,
nicht Deployment, Signing oder Fach- und Sicherheitsanforderungen
\parencite{kleiveist2026acceptance}. Das Produkt benötigt dafür eigene Tests
und Nachweise.
Diese müssen die konkreten Kundenanforderungen, Integrationen und
Betriebsbedingungen abdecken.
\beginnerinput{workspace/statement/800_einsatz_transfer/860}
